% Original timeline; winding path inspired by user reference, with no decorative icons.
\begin{tikzpicture}[x=1mm,y=1mm,font=\figurefont\footnotesize,
 milestone/.style={circle,minimum size=7mm,inner sep=0pt,draw=BookTeal,fill=BookTeal!20,line width=.8pt},
 title/.style={font=\figurefont\small,text=BookInk},
 detail/.style={font=\figurefont\footnotesize,text=BookMuted}]
\path[use as bounding box] (0,-8) rectangle (155,182);
\draw[BookMuted!60,line width=3.5mm,line cap=round]
 (78,178) .. controls (48,174) and (48,164) .. (78,158)
 .. controls (108,152) and (108,142) .. (78,136)
 .. controls (48,130) and (48,120) .. (78,114)
 .. controls (108,108) and (108,98) .. (78,92)
 .. controls (48,86) and (48,76) .. (78,70)
 .. controls (108,64) and (108,54) .. (78,48)
 .. controls (48,42) and (48,32) .. (78,26)
 .. controls (108,20) and (108,10) .. (78,4);
\draw[white,dashed,line width=.5mm,dash pattern=on 2mm off 2mm]
 (78,178) .. controls (48,174) and (48,164) .. (78,158)
 .. controls (108,152) and (108,142) .. (78,136)
 .. controls (48,130) and (48,120) .. (78,114)
 .. controls (108,108) and (108,98) .. (78,92)
 .. controls (48,86) and (48,76) .. (78,70)
 .. controls (108,64) and (108,54) .. (78,48)
 .. controls (48,42) and (48,32) .. (78,26)
 .. controls (108,20) and (108,10) .. (78,4);
\node[milestone] at (56,170) {};
\node[anchor=west,text=BookTeal,font=\figurefont\small] at (2,175) {1805—1809};
\node[title,anchor=west,text width=49mm,align=left] at (2,169) {最小二乘};
\node[detail,anchor=north west,text width=49mm,align=left] at (2,165) {用平方残差组织参数拟合};
\node[milestone] at (100,148) {};
\node[anchor=west,text=BookTeal,font=\figurefont\small] at (109,153) {1901};
\node[title,anchor=west,text width=46mm,align=left] at (109,147) {主成分分析};
\node[detail,anchor=north west,text width=46mm,align=left] at (109,143) {用线性子空间压缩数据};
\node[milestone] at (56,126) {};
\node[anchor=west,text=BookTeal,font=\figurefont\small] at (2,131) {1936};
\node[title,anchor=west,text width=49mm,align=left] at (2,125) {线性判别};
\node[detail,anchor=north west,text width=49mm,align=left] at (2,121) {用类别统计量构造分界};
\node[milestone] at (100,104) {};
\node[anchor=west,text=BookTeal,font=\figurefont\small] at (109,109) {1967};
\node[title,anchor=west,text width=46mm,align=left] at (109,103) {近邻方法};
\node[detail,anchor=north west,text width=46mm,align=left] at (109,99) {按样本距离形成局部预测};
\node[milestone] at (56,82) {};
\node[anchor=west,text=BookTeal,font=\figurefont\small] at (2,87) {1984};
\node[title,anchor=west,text width=49mm,align=left] at (2,81) {分类与回归树};
\node[detail,anchor=north west,text width=49mm,align=left] at (2,77) {用递归划分组织局部规则};
\node[milestone] at (100,60) {};
\node[anchor=west,text=BookTeal,font=\figurefont\small] at (109,65) {1995};
\node[title,anchor=west,text width=46mm,align=left] at (109,59) {支持向量机};
\node[detail,anchor=north west,text width=46mm,align=left] at (109,55) {用间隔与核映射构造分界};
\node[milestone] at (56,38) {};
\node[anchor=west,text=BookTeal,font=\figurefont\small] at (2,43) {1996};
\node[title,anchor=west,text width=49mm,align=left] at (2,37) {Bagging};
\node[detail,anchor=north west,text width=49mm,align=left] at (2,33) {组合扰动样本上学得的模型};
\node[milestone] at (100,16) {};
\node[anchor=west,text=BookTeal,font=\figurefont\small] at (109,21) {2001};
\node[title,anchor=west,text width=46mm,align=left] at (109,15) {随机森林与梯度提升};
\node[detail,anchor=north west,text width=46mm,align=left] at (109,11) {以随机化或残差拟合组织集成};
\end{tikzpicture}
